\chapter{\IfLanguageName{dutch}{Stand van zaken}{State of the art}}%
\label{ch:stand-van-zaken}

\subsection{Het belang van veilig wachtwoordbeheer}

Het belang van veilig wachtwoordbeheer wordt in de literatuur vaak benadrukt. Gebruikers die hun wachtwoorden hergebruiken over meerdere sites lopen een verhoogd risico op cyberaanvallen \autocite{Gaw2006}. Studies tonen bovendien aan dat het gebruik van wachtwoordmanagers niet langer een luxe is, maar een noodzaak voor het behouden van digitale veiligheid \autocite{CMU2025}.

\subsection{Typen wachtwoordmanagers}
\textcite{Zawalnyski2025} verdeelt wachtwoordmanagers in drie categorieën: browser, cloud en desktopgebaseerd. Elk type heeft zijn voor- en nadelen, afhankelijk van de behoeften van de gebruiker.

Browsergebaseerde wachtwoordmanagers zijn populair omdat ze gratis en eenvoudig te gebruiken zijn \autocite{Rimasiute2025}. Verschillende internetbrowsers, zoals Google Chrome, Firefox, Edge en Safari, bieden ingebouwde wachtwoordmanagers aan waarmee gebruikers wachtwoorden kunnen opslaan en automatisch laten invullen. Een nadeel is dat wachtwoorden niet altijd eenvoudig kunnen worden gesynchroniseerd tussen verschillende browsers. Daarnaast kan het onveilig zijn om dit type wachtwoordmanager te gebruiken op gedeelde apparaten, waar meerdere gebruikers toegang hebben tot dezelfde browser \autocite{Zoho}.
Volgens \textcite{Pearman2019} gebruiken gebruikers van browsergebaseerde wachtwoordmanagers vaker zwakke wachtwoorden en hergebruiken zij deze ook vaker. De interviews uit hun studie suggereren dat deze gebruikers vooral gedreven worden door gebruiksgemak in plaats van door veiligheid.

Cloudgebaseerde wachtwoordmanagers worden meestal gedownload en beheerd op mobiele apparaten of desktops. Dit type slaat wachtwoorden op in een beveiligde cloudomgeving. In tegenstelling tot browsergebaseerde wachtwoordmanagers kunnen cloudgebaseerde zwakke en hergebruikte wachtwoorden detecteren. Verder voorzien de meeste ook een back-up van de wachtwoordkluis en werken ze op meerdere browsers en besturingssystemen \autocite{Rimasiute2025}. Dit kan het gebruik veiliger maken dan bij browsergebaseerde wachtwoordmanagers. Een nadeel is dat gebruikers minder controle hebben over de beveiliging van de wachtwoordkluis, aangezien hun gegevens op servers van een derde partij servers worden opgeslagen.

Desktop-gebaseerde wachtwoordmanagers worden lokaal op de computer geïnstalleerd en slaan wachtwoorden lokaal op. Individuen en bedrijven die een wachtwoordmanager in eigen omgeving willen hosten en beheren, geven doorgaans de voorkeur geven aan dit type \autocite{Zoho}. hierbij is geen derde partij betrokken. Gebruikers blijven wel verantwoordelijk voor het maken van handmatige back-ups, voor het geval het toestel crasht. Verder is het noodzakelijk dat gebruikers hun computer goed beveiligen en een sterk masterwachtwoord gebruiken om ongewenste toegang te voorkomen.

\subsection{Drempels voor beginnende gebruikers}
Ondanks de voordelen blijkt uit onderzoek dat veel beginnende gebruikers terughoudend zijn om wachtwoordmanagers te gebruiken. Volgens \textcite{Pearman2019} hebben gebruikers vaak onvoldoende vertrouwen in de beveiliging van deze tools en ervaren de werking ervan als complex. Daarnaast weten sommige gebruikers niet dat er betere opties bestaan dan het simpelweg onthouden van wachtwoorden of het opschrijven ervan.  
Uit onderzoek van \textcite{Fagan2017} blijkt dat 45,9\% van de niet-gebruikers aangeeft zich zorgen te maken over de veiligheid van wachtwoordmanagers. Ze geven aan bang te zijn dat een wachtwoordmanager gecompromitteerd kan worden en vinden het daarom risicovol om al hun wachtwoorden op één centrale locatie op te slaan. Een andere 42,3\% van de niet-gebruikers geeft aan geen nood te hebben aan het gebruik van een wachtwoordmanager. Dit kan wijzen op een beperkte kennis van de veiligheidsrisico's en/of voordelen van dergelijke tools.
Ook \textcite{Ray2021} geeft aan dat een gebrek aan vertrouwen een belangrijke drempel vormt en dat niet-gebruikers vaak de voorkeur geven aan het zelf beheren van hun wachtwoorden, in plaats van dit toe te vertrouwen aan een wachtwoordmanager. 

\subsection{Beperkingen van traditionele leermethoden}
Binnen security awareness wordt steeds duidelijker dat traditionele methoden, zoals tekstuele handleidingen of tutorials, niet voldoende zijn om gebruikers effectief te trainen in veilig wachtwoordbeheer. Een belangrijke oorzaak is dat gebruikers wel advies krijgen over veilig gedrag, maar dit niet in de praktijk kunnen brengen zonder risico op fouten of datalekken \autocite{Sugatan2020}. 

\subsection{Interactieve leeromgevingen als oplossing}
\textcite{Sugatan2020} gebruikten in hun onderzoek een interactieve, verhalende leeromgeving waar gebruikers de kans kregen om beslissingen te nemen rond wachtwoordgedrag en vervolgens ook de gevolgen te zien. De onderzoekers stelden vast dat deze aanpak leidde tot een beter begrip, meer vertrouwen en een grotere bereidheid om zelf een wachtwoordmanager te gebruiken.
 Dit onderzoek onderstreept het potentieel van interactieve simulaties als een effectieve methode om beginnende gebruikers te trainen in veilig wachtwoordbeheer.
























% Tip: Begin elk hoofdstuk met een paragraaf inleiding die beschrijft hoe
% dit hoofdstuk past binnen het geheel van de bachelorproef. Geef in het
% bijzonder aan wat de link is met het vorige en volgende hoofdstuk.

% Pas na deze inleidende paragraaf komt de eerste sectiehoofding.

% Dit hoofdstuk bevat je literatuurstudie. De inhoud gaat verder op de inleiding, maar zal het onderwerp van de bachelorproef *diepgaand* uitspitten. De bedoeling is dat de lezer na lezing van dit hoofdstuk helemaal op de hoogte is van de huidige stand van zaken (state-of-the-art) in het onderzoeksdomein. Iemand die niet vertrouwd is met het onderwerp, weet nu voldoende om de rest van het verhaal te kunnen volgen, zonder dat die er nog andere informatie moet over opzoeken \autocite{Pollefliet2011}.

% Je verwijst bij elke bewering die je doet, vakterm die je introduceert, enz.\ naar je bronnen. In \LaTeX{} kan dat met het commando \texttt{$\backslash${textcite\{\}}} of \texttt{$\backslash${autocite\{\}}}. Als argument van het commando geef je de ``sleutel'' van een ``record'' in een bibliografische databank in het Bib\LaTeX{}-formaat (een tekstbestand). Als je expliciet naar de auteur verwijst in de zin (narratieve referentie), gebruik je \texttt{$\backslash${}textcite\{\}}. Soms is de auteursnaam niet expliciet een onderdeel van de zin, dan gebruik je \texttt{$\backslash${}autocite\{\}} (referentie tussen haakjes). Dit gebruik je bv.~bij een citaat, of om in het bijschrift van een overgenomen afbeelding, broncode, tabel, enz. te verwijzen naar de bron. In de volgende paragraaf een voorbeeld van elk.

% \textcite{Knuth1998} schreef een van de standaardwerken over sorteer- en zoekalgoritmen. Experten zijn het erover eens dat cloud computing een interessante opportuniteit vormen, zowel voor gebruikers als voor dienstverleners op vlak van informatietechnologie~\autocite{Creeger2009}.

% Let er ook op: het \texttt{cite}-commando voor de punt, dus binnen de zin. Je verwijst meteen naar een bron in de eerste zin die erop gebaseerd is, dus niet pas op het einde van een paragraaf.

% \begin{figure}
%   \centering
%   \includegraphics[width=0.8\textwidth]{grail.jpg}
%   \caption[Voorbeeld figuur.]{\label{fig:grail}Voorbeeld van invoegen van een figuur. Zorg altijd voor een uitgebreid bijschrift dat de figuur volledig beschrijft zonder in de tekst te moeten gaan zoeken. Vergeet ook je bronvermelding niet!}
% \end{figure}

% \begin{listing}
%   \begin{minted}{python}
%     import pandas as pd
%     import seaborn as sns

%     penguins = sns.load_dataset('penguins')
%     sns.relplot(data=penguins, x="flipper_length_mm", y="bill_length_mm", hue="species")
%   \end{minted}
%   \caption[Voorbeeld codefragment]{Voorbeeld van het invoegen van een codefragment.}
% \end{listing}

% \lipsum[7-20]

% \begin{table}
%   \centering
%   \begin{tabular}{lcr}
%     \toprule
%     \textbf{Kolom 1} & \textbf{Kolom 2} & \textbf{Kolom 3} \\
%     $\alpha$         & $\beta$          & $\gamma$         \\
%     \midrule
%     A                & 10.230           & a                \\
%     B                & 45.678           & b                \\
%     C                & 99.987           & c                \\
%     \bottomrule
%   \end{tabular}
%   \caption[Voorbeeld tabel]{\label{tab:example}Voorbeeld van een tabel.}
% \end{table}

